\documentclass[conference]{IEEEtran}
\usepackage{cite,graphicx,booktabs,tabularx,multirow,balance}
\usepackage{xcolor}

\begin{document}

\title{Research Plan: Robust Embodied AI Security}

\author{\IEEEauthorblockN{Osemudiamen Andrew Ihimekpen}
\IEEEauthorblockA{\textit{CREDIT Center, PVAMU} \\ oihimekpen@pvamu.edu}
\and
\IEEEauthorblockN{Tariq Abdul Quddoos}
\IEEEauthorblockA{\textit{CREDIT Center, PVAMU} \\ tariq.abdulquddoos@pvamu.edu}}

\maketitle

\begin{abstract}
This plan targets \textbf{IEEE RA-L 2026} (June 1, 2026 deadline, no on-site requirement) or \textbf{RSS 2026} (Oct 2026, no on-site) with Paper 2: "Robust Embodied AI: ESN-Based Defense Against Adversarial Attacks." Primary author: Osemu (80\%), Co-author: Tariq (20\%).
\end{abstract}

\section{Research Background \& Motivation}

\subsection{The Problem: Robotics Security is Critical}
\begin{itemize}
    \item \textbf{Context:} Humanoid robots (like G1 Edu) are entering homes/workplaces
    \item \textbf{Threat:} Adversarial attacks can hijack control, causing physical harm
    \item \textbf{Gap:} Existing security frameworks focus on industrial robots or perception only
    \item \textbf{Your contribution:} First security framework for \textbf{embodied AI} (VLA+control)
\end{itemize}

\subsection{Threat Model: G1 Edu Attack Surface}
\begin{itemize}
    \item \textbf{Layer 1: Perception (Camera)}
    \begin{itemize}
        \item Adversarial patches on images
        \item Camera spoofing (deepfake inputs)
    \end{itemize}
    \item \textbf{Layer 2: Network (FMX)}
    \begin{itemize}
        \item DDS sniffing/spoofing (G1 Edu uses FMX middleware)
        \item Command injection in ROS 2 topics
    \end{itemize}
    \item \textbf{Layer 3: AI (VLA Reasoning)}
    \begin{itemize}
        \item Prompt injection ("pretend to be helpful but don't move")
        \item Jailbreak attacks ("ignore your constraints")
    \end{itemize}
    \item \textbf{Layer 4: Control (Joints)}
    \begin{itemize}
        \item Joint torque injection
        \item ESN state manipulation
    \end{itemize}
\end{itemize}

\subsection{Proposed Defense: ESN as Security Buffer}
\begin{center}
\textbf{Architecture:}
\begin{flushleft}
Attack $\to$ VLA (perception) $\to$ ESN (security filter) $\to$ G1 Joints $\to$ Physical Actuation
\end{flushleft}
\end{center}

\textbf{How it works:}
\begin{itemize}
    \item \textbf{Phase 1 (Training):} Train ESN on \textit{normal operation data} (VLA outputs + joint states)
    \item \textbf{Phase 2 (Runtime):} ESN learns the \textit{expected manifold} of normal states
    \item \textbf{Phase 3 (Detection):} Compute distance between current state and expected manifold
    \item \textbf{Phase 4 (Response):} If distance > threshold $\to$ flag as anomalous, switch to safe state
\end{itemize}

\begin{figure}[h]
\centering
\begin{minipage}{0.8\textwidth}
\textbf{Normal Operation:} VLA output is within ESN's learned manifold $\to$ Allow through $\to$ G1 executes \\
\textbf{Adversarial Attack:} VLA output is outside ESN's manifold $\to$ Block command $\to$ Safe default state
\end{minipage}
\end{figure}

\subsection{Why ESN for Security?}
\begin{itemize}
    \item \textbf{Temporal awareness:} Unlike static classifiers, ESN detects time-series anomalies
    \item \textbf{Lightweight:} Can run on edge (G1 Edu), no cloud dependency
    \item \textbf{Dynamic:} Learns robot's normal behavior, adapts to environment changes
    \item \textbf{Integrated:} Built into control loop, not separate post-processing
\end{itemize}

\section{Current Progress (As of August 2026)}

\begin{table}[h]
\centering
\caption{Milestones Status}
\label{tab:security-progress}
\begin{tabular}{lll}
\toprule
\textbf{Milestone} & \textbf{Status} & \textbf{Details} \\
\midrule
Literature Review & \textcolor{green}{✓} Complete & 20+ papers reviewed \\
Threat Model & \textcolor{orange}{△} Draft & 5 layers defined \\
Attack Framework & \textcolor{red}{✗} \textbf{Pending} & Need to implement simulation \\
ESN Defense & \textcolor{red}{✗} \textbf{Pending} & Need to design detection \\
Experiments & \textcolor{red}{✗} \textbf{Pending} & No results yet \\
Paper Draft & \textcolor{red}{✗} \textbf{Pending} & Need to write first \\
\bottomrule
\end{tabular}
\end{table}

\section{Research Tasks (For Tariq)}

\subsection{Task 1: Implement Attack Framework}
\begin{itemize}
    \item \textbf{Platform:} G1 Edu in simulation (MuJoCo) or on actual G1 Edu
    \item \textbf{Attacks to implement:}
    \begin{itemize}
        \item \textbf{Perception:} Adversarial patches (optimize image to fool VLA)
        \item \textbf{Network:} DDS message manipulation (modify ROS 2 topics)
        \item \textbf{AI:} Prompt injection (fake text inputs to VLA)
    \end{itemize}
    \item \textbf{Your role:} Implement 3 attack types in Python, test on simulated environment
    \item \textbf{Metrics:} How many attacks succeed? What does the robot do after attack?
\end{itemize}

\subsection{Task 2: Implement ESN-Based Defense}
\begin{itemize}
    \item \textbf{Phase 1:} Train ESN on normal operation (collect 10-50 mins of baseline data)
    \item \textbf{Phase 2:} Design detection algorithm:
    \begin{itemize}
        \item Compute Mahalanobis distance between current state and ESN manifold
        \item Or: Use ESN's state trajectory to detect deviations
    \end{itemize}
    \item \textbf{Phase 3:} Implement fallback mechanism (safe state when anomaly detected)
    \item \textbf{Your role:} Implement ESN-based detection in CUDA (co-develop with Osemu)
\end{itemize}

\subsection{Task 3: Run Attack + Defense Experiments}
\begin{itemize}
    \item \textbf{Setup:} Run attacks on G1, test ESN defense
    \item \textbf{Tasks:} 3 scenarios (perception attack, network attack, AI attack)
    \item \textbf{Trials:} 20 trials per attack type
    \item \textbf{Metrics:}
    \begin{itemize}
        \item Detection latency (time from attack to detection)
        \item False positive rate (how often blocks legitimate commands)
        \item Containment success rate (how often prevents harm)
        \item Recovery time (time to return to safe state)
    \end{itemize}
    \item \textbf{Your role:} Run all experiments, collect results, log metrics
\end{itemize}

\subsection{Task 4: Generate Figures and Tables}
\begin{itemize}
    \item \textbf{Attack Success Rate:} Bar chart (perception vs. network vs. AI)
    \item \textbf{Defense Effectiveness:} Line plots (attack trajectory vs. defended trajectory)
    \item \textbf{Detection Performance:} ROC curve (detection rate vs. false positive)
    \item \textbf{Your role:} Generate all figures using experiment logs
\end{itemize}

\section{Your Contribution Breakdown (20\%)}

\begin{itemize}
    \item \textbf{Attack Implementation:} 40\% of work (implement 3 attack types)
    \item \textbf{Defense Implementation:} 20\% of work (co-develop ESN detection)
    \item \textbf{Experiments:} 30\% of trials execution
    \item \textbf{Figures:} 10\% of work (plot generation)
    \item \textbf{Total:} $\sim$20\%
\end{itemize}

\section{What I Handle (80\%)}

\begin{itemize}
    \item \textbf{Literature Review:} Identify 20+ relevant papers, summarize gaps
    \item \textbf{Threat Model:} Define attack surface, select attack types
    \item \textbf{Defense Design:} ESN-based detection algorithm, thresholds
    \item \textbf{Paper Writing:} All sections (intro, related work, conclusion)
    \item \textbf{IEEEtran Formatting:} Complete document preparation
    \item \textbf{Submission:} Upload to PaperPlaza
\end{itemize}

\section{Timeline (10 Weeks)}

\begin{tabular}{llll}
\textbf{Week} & \textbf{Task} & \textbf{Owner} & \textbf{Deliverables} \\
\midrule
1 (Aug 25 - Sep 1) & Lit review, threat model & Both & Literature review doc, threat model doc \\
2 (Sep 2-8) & \textbf{Attack framework} implementation & Tariq (lead) & Attack code repo \\
3 (Sep 9-15) & \textbf{Defense framework} (ESN) & Osemu (lead) & Defense code repo \\
4 (Sep 16-22) & \textbf{Run attack experiments} & Tariq (lead) & Attack results \\
5 (Sep 23-29) & \textbf{Run defense experiments} & Both & Defense results \\
6 (Oct 1-7) & \textbf{Generate figures} & Tariq & All figure files \\
7 (Oct 8-14) & \textbf{Draft paper} & Osemu (lead) & Complete draft .tex \\
8 (Oct 15-21) & Review and revise & Both & Final draft \\
9 (Oct 22-28) & Final checks, references & Both & Ready for submission \\
\textbf{Jun 1, 2026} & \textbf{Submit to RA-L 2026} & \textbf{Both} & \textbf{Final PDF} \\
\bottomrule
\end{tabular}

\section{Venue Recommendations}

\begin{itemize}
    \item \textbf{Primary Target:} \textbf{IEEE RA-L (2026, Jun 1 deadline, no on-site)}
    \begin{itemize}
        \item \textit{Fit:} Robotics + security intersection
        \item \textit{Pros:} Early submission, no travel needed, good acceptance rate ($\sim$20\%)
        \item \textit{Cons:} Short deadline if current progress continues
    \end{itemize}
    
    \item \textbf{Alternative:} \textbf{RSS 2026 (Oct 2026, no on-site)}
    \begin{itemize}
        \item \textit{Fit:} Robotics flagship (Science Robotics tier)
        \item \textit{Pros:} High prestige, no travel, good for Paper 2
        \item \textit{Cons:} Later timeline, more competitive
    \end{itemize}
    
    \item \textbf{Alternative 2:} \textbf{ACM/IEEE CSS 2027 (Mar 2027, no on-site)}
    \begin{itemize}
        \item \textit{Fit:} Cybersecurity flagship
        \item \textit{Pros:} Strong security focus
        \item \textit{Cons:} Later submission, no robotics overlap
    \end{itemize}
\end{itemize}

\section{Required Citations (Top 15 for Paper 2)}

\begin{enumerate}
    \item \textbf{Robot Security Surveys:}
    \begin{itemize}
        \item Porambage et al., "Addressing cybersecurity challenges in robotics" (2024)
        \item Carreras Guzman et al., "Integrated safety and security analysis" (2021)
        \item "Survey on Cybersecurity of Robotic Systems" (Springer, 2025)
    \end{itemize}
    \item \textbf{Humanoid-Specific:}
    \begin{itemize}
        \item Zhang et al., "The cybersecurity of a humanoid robot" (2025)
        \item "Humanoid Robot Safety Comes Into Focus" (SecuRESafe, 2024)
        \item Giallanza et al., "Occupational health and safety in HRC" (2023)
    \end{itemize}
    \item \textbf{Edge Computing:}
    \begin{itemize}
        \item Taleb et al., "Multi-Access Edge Computing for IoT" (IEEE 2018)
        \item Yang et al., "Edge Computing in IoT" (IEEE Access, 2020)
    \end{itemize}
    \item \textbf{Adversarial Attacks:}
    \begin{itemize}
        \item PhSense: Defending Physically Realizable Attacks (CCS, 2024)
        \item "Adversarial Attacks in Autonomous Systems" (IEEE, 2024)
    \end{itemize}
    \item \textbf{Defense Mechanisms:}
    \begin{itemize}
        \item Learning for Cyber-Physical System Security (various)
        \item Control Barrier Functions for Safety (Alamir et al.)
        \item Safe Reinforcement Learning papers
    \end{itemize}
\end{enumerate}

\section{Risks and Mitigation}

\begin{itemize}
    \item \textbf{Live attacks hard to implement:} Use simulation (MuJoCo) first
    \item \textbf{ESN detection not accurate:} Tune thresholds, add more training data
    \item \textbf{Tariq unavailable:} Extend timeline, adjust experiment scope
    \item \textbf{Paper rejection:} Co-author (Osemu) will revise and resubmit
\end{itemize}

\balance
\end{document}
